\documentclass[11pt]{article}
\usepackage[T1]{fontenc}
\usepackage[utf8]{inputenc}
\usepackage{amsmath,amsthm,mathtools}
\usepackage{newtxtext,newtxmath}
\usepackage[margin=1in,headheight=15pt]{geometry}
\usepackage{microtype}
\usepackage{booktabs,array,longtable}
\usepackage[dvipsnames]{xcolor}
\usepackage{enumitem}
\usepackage{needspace}
\usepackage{fancyhdr}
\usepackage{titlesec}
\usepackage{listings}
\usepackage[most]{tcolorbox}
\usepackage{hyperref}
\usepackage[nameinlink,noabbrev]{cleveref}
\definecolor{Ink}{HTML}{17324D}
\definecolor{Accent}{HTML}{126D76}
\definecolor{Pale}{HTML}{F0F6F7}
\hypersetup{colorlinks=true,linkcolor=Ink,citecolor=Accent,urlcolor=Accent,
 pdftitle={Binary Carries and the Divisibility of Two-Stack-Sortable Permutation Numbers},
 pdfauthor={Research article prepared with ChatGPT}}
\titleformat{\section}{\Large\bfseries\color{Ink}}{\thesection}{0.7em}{}
\titleformat{\subsection}{\large\bfseries\color{Accent}}{\thesubsection}{0.7em}{}
\pagestyle{fancy}\fancyhf{}
\fancyhead[L]{\small\color{Ink}Binary carries and A000139}
\fancyhead[R]{\small Research article}
\fancyfoot[C]{\small\thepage}
\renewcommand{\headrulewidth}{0.3pt}
\newtheorem{theorem}{Theorem}[section]
\newtheorem{lemma}[theorem]{Lemma}
\newtheorem{proposition}[theorem]{Proposition}
\newtheorem{corollary}[theorem]{Corollary}
\theoremstyle{definition}\newtheorem{definition}[theorem]{Definition}
\theoremstyle{remark}\newtheorem{remark}[theorem]{Remark}
\newcommand{\N}{\mathbb N}
\newcommand{\Z}{\mathbb Z}
\newcommand{\Q}{\mathbb Q}
\newcommand{\F}{\mathbb F}
\newcommand{\E}{\mathbb E}
\newcommand{\Prob}{\mathbb P}
\newcommand{\Var}{\operatorname{Var}}
\newcommand{\bit}{\operatorname{s}_2}
\newcommand{\val}{\nu_2}
\newcommand{\bits}{\operatorname{bits}}
\newcommand{\coeff}[2]{[#1^{#2}]}
\newcommand{\ph}{\varphi}
\newcommand{\dd}{\mathop{}\!\mathrm d}
\newcommand{\ind}{\mathbf 1}
\newcommand{\vseq}{v}
\newcommand{\rank}{\operatorname{rank}}
\newcommand{\rhoo}{\rho}
\lstset{basicstyle=\small\ttfamily,breaklines=true,columns=fullflexible,
 frame=single,rulecolor=\color{Ink!25},backgroundcolor=\color{Pale},
 showstringspaces=false,keepspaces=true}
\setlist{nosep,leftmargin=1.6em}
\setlength{\parskip}{4pt}
\setlength{\parindent}{1.2em}
\emergencystretch=1.5em
\allowdisplaybreaks[2]

\begin{document}
\begin{center}
{\small\sffamily\color{Accent}COMBINATORICS \quad INTEGER SEQUENCES \quad ASYMPTOTICS}\par\vspace{0.7cm}
{\LARGE\bfseries\color{Ink}Binary Carries and the Divisibility of\\[0.15cm]
Two-Stack-Sortable Permutation Numbers}\par\vspace{0.45cm}
{\large A proof of the OEIS A000139 parity conjecture,\\
with exact valuation distributions and limit laws}\par\vspace{0.5cm}
{September 19, 2026}
\end{center}
\vspace{0.2cm}
\begin{abstract}
For $n\geq 0$, put
\[
 a_n=\frac{2(3n)!}{(n+1)!(2n+1)!},\qquad v(n)=\nu_2(a_n).
\]
The OEIS entry A000139 records a July 2025 conjecture of Peter Bala:
$a_n$ is odd exactly when $n$ is an odd fibbinary number. We prove the
assertion in two independent ways, using binary carries and formal power
series over $\F_2$. More generally, we determine the rational bivariate
generating function of the complete valuation distribution on every
interval $0\leq n<2^m$. It has a cubic denominator, giving a third-order
recurrence in the binary length. For each fixed valuation $r$, the counts
have an explicitly determined minimal eventual recurrence and asymptotics
of order $m^{\lfloor r/2\rfloor}\ph^m$, with explicit constants. We also
classify maximal valuations and find the least positive index of every
valuation. For a uniformly chosen index below $2^m$, we obtain exact mean
and variance formulas, a central limit theorem, and a uniform local limit
theorem. A six-state weighted digit automaton computes the full distribution
below an arbitrary bound without enumerating its indices. The accompanying
Python programs, exact data, and execution records reproduce the finite
checks. The proofs do not depend on those checks.
\end{abstract}

\begin{tcolorbox}[colback=Pale,colframe=Accent,boxrule=0.6pt,arc=1mm]
\textbf{Main resolution.} For every $n\geq0$,
\[
 \boxed{\ a_n\text{ is odd}\quad\Longleftrightarrow\quad
 n\text{ is odd and its binary expansion contains no }11.\ }
\]
\textbf{Status of the claim.} The assertion was still labeled a conjecture
in the OEIS entry inspected for this article. This is a proof of that
recorded assertion, not a claim that an exhaustive literature search has
established first priority. The valuation refinements are proved here as
well; their priority has not been independently established.
\end{tcolorbox}
\newpage
{\small\setlength{\parskip}{0pt}\tableofcontents}
\newpage

\section{The selected problem and the scope of the result}\label{sec:setup}

The integer sequence
\begin{equation}\label{eq:a}
 a_n=\frac{2(3n)!}{(n+1)!(2n+1)!}
 =\frac{2}{(n+1)(2n+1)}\binom{3n}{n}
\end{equation}
begins
\[
 2,1,2,6,22,91,408,1938,9614,49335,260130,1402440,\ldots.
\]
For $n\geq1$ it counts permutations sorted by two successive applications
of the usual stack-sorting map. This is the deterministic stack-sorting
operator, not an arbitrary two-stack machine with unrestricted choices.
The same formula occurs in the enumeration of fighting fish; Fang gives a
bijection between these objects and two-stack-sortable permutations
\cite{Fang}. The formula itself and several further interpretations are
recorded in OEIS A000139 \cite{OEIS139}. Our arithmetic arguments use only
\eqref{eq:a}.

An odd \emph{fibbinary number} is an odd nonnegative integer whose binary
expansion contains no two consecutive $1$s. These numbers are
\[
 1,5,9,17,21,33,37,41,65,69,73,81,85,\ldots,
\]
and form OEIS A022341 \cite{OEIS22341}. The terminology and its connection
with Fibonacci representations are discussed by Lindroos, Sills, and Wang
\cite{LSW}.

The selected question is the assertion attributed to Peter Bala, dated
July 24, 2025, in A000139: the indices of the odd terms are exactly these
odd fibbinary numbers \cite{OEIS139}. The entry inspected on September 19,
2026 still used the conjecture label. Its inspected internal version
identified itself as revision 297, dated May 12, 2026 \cite{OEIS139internal}.
That documents the provenance of the question. It does not exclude the
possibility of an earlier proof elsewhere or of an unrecorded resolution.

The problem is well suited to an elementary attack because a factorial
quotient converts prime divisibility into digit sums. At the prime $2$,
the relevant sum is $n+2n$, so the two summands differ by just a binary
shift. The apparent global arithmetic question becomes a finite-state
question about adjacent digits and propagated carries. We pursue this
structure rather than stopping at a parity test.

\subsection{Conventions}
For a positive integer $b$, $\nu_2(b)$ denotes the exponent of $2$ in $b$.
We write $\bit(n)$ for the number of $1$s in the binary expansion of $n$,
with $\bit(0)=0$. Define
\[
 t(n)=\nu_2(n+1),\qquad
 c(n)=\text{the total number of carries in the binary addition }n+2n.
\]
Every outgoing carry is counted, including a carry beyond the highest
nonzero input digit. The number $t(n)$ is the length of the terminal run
of $1$s in the usual, most-significant-digit-first binary expansion.

The index $n=0$ is retained throughout: $a_0=2$ and $v(0)=1$. Thus it
contributes one term of valuation $1$, not one odd term. A dyadic interval
always means $0\leq n<2^m$, and an arbitrary bound is also exclusive.
Fibonacci numbers are normalized by $F_0=0$, $F_1=1$. Finally,
\[
 \ph=\frac{1+\sqrt5}{2},\qquad \rho=\ph^{-1},\qquad
 \alpha=\log_2\ph.
\]

\Needspace{8\baselineskip}
\section{An elementary carry proof of the parity conjecture}\label{sec:parity}

\subsection{Factorial valuations and digit sums}
\begin{lemma}\label{lem:legendre}
For every integer $k\geq0$,
\[
 \nu_2(k!)=k-\bit(k).
\]
\end{lemma}
\begin{proof}
Each multiple of $2^j$ contributes one additional factor of $2$ to $k!$,
so its valuation is $\sum_{j\geq1}\lfloor k/2^j\rfloor$. Write
$k=\sum_{i\geq0}b_i2^i$, where $b_i\in\{0,1\}$. Interchanging finite sums,
\[
 \sum_{j\geq1}\left\lfloor\frac{k}{2^j}\right\rfloor
 =\sum_{i\geq1}b_i\sum_{j=1}^i2^{i-j}
 =\sum_{i\geq0}b_i(2^i-1)=k-\bit(k).
\]
\end{proof}

\begin{proposition}[Exact valuation]\label{prop:valuation}
For every $n\geq0$,
\begin{align}
 v(n)&=\bit(n)+\bit(n+1)-\bit(3n),\label{eq:digitval}\\
     &=1-t(n)+c(n).\label{eq:carryval}
\end{align}
\end{proposition}
\begin{proof}
Apply \cref{lem:legendre} to \eqref{eq:a}. The result is
\[
 1+(3n-\bit(3n))-(n+1-\bit(n+1))
 -(2n+1-\bit(2n+1)).
\]
Since $\bit(2n+1)=\bit(n)+1$, this is \eqref{eq:digitval}.
Incrementing $n$ changes its $t(n)$ terminal $1$s into $0$s and changes
the next $0$ into a $1$. Consequently,
\begin{equation}\label{eq:trailing}
 \bit(n+1)=\bit(n)+1-t(n).
\end{equation}

For completeness, the carry identity needed here follows column by
column. Let $b_i$ be the digits of $n$, put $b_{-1}=0$, and let $q_i$ be
the incoming carry in column $i$, with $q_0=0$. If $d_i$ is the output
digit of $3n$, then
\[
 b_i+b_{i-1}+q_i=d_i+2q_{i+1}.
\]
Sum this identity over all columns, including zero columns until the
carry has vanished. The sums of the incoming and outgoing carries agree,
and therefore
\begin{equation}\label{eq:carrydigits}
 c(n)=\sum_{i\geq0}q_{i+1}=2\bit(n)-\bit(3n).
\end{equation}
Combining \eqref{eq:trailing} and \eqref{eq:carrydigits} gives
\eqref{eq:carryval}.
\end{proof}

\subsection{Why the trailing run matters}
It would be incorrect to infer parity solely from the absence of carries
in $\binom{3n}{n}$. The factor $2/(n+1)$ in \eqref{eq:a} changes the
valuation by $1-t(n)$. The following three cases account for this
correction exactly.

\begin{theorem}[The recorded parity conjecture]\label{thm:parity}
For every $n\geq0$, the following are equivalent:
\begin{enumerate}[label=(\roman*)]
\item $a_n$ is odd;
\item $n$ is odd and the binary expansion of $n$ contains no adjacent $1$s;
\item $n$ is odd and the bitwise expression $n\mathbin{\&}(2n)$ is zero.
\end{enumerate}
Thus the odd indices are precisely OEIS A022341.
\end{theorem}
\begin{proof}
If $n$ is even, then $t(n)=0$, so \eqref{eq:carryval} gives
$v(n)=1+c(n)\geq1$.

Suppose $n$ ends in exactly one $1$, so $t(n)=1$. Then $v(n)=c(n)$.
There is no carry in $n+2n$ exactly when the nonzero digits of the two
summands do not overlap. Indeed, an overlap creates a carry, while in
the absence of overlaps the incoming carry remains zero inductively.
Because the second summand is a shift of the first, an overlap is exactly
an adjacent pair $11$ in $n$. Hence $v(n)=0$ exactly in the desired case.

It remains to exclude a terminal run of length $t\geq2$. Its
digits satisfy $b_0=\cdots=b_{t-1}=1$ and $b_t=0$. In column $1$,
$b_1+b_0=2$ starts a carry. That carry continues through columns
$2,\ldots,t-1$. In column $t$ we still have
$b_t+b_{t-1}+q_t=0+1+1=2$. There are consequently at least $t$ outgoing
carries, in columns $1,\ldots,t$. Equation \eqref{eq:carryval} now yields
$v(n)\geq1-t+t=1$. This excludes every remaining odd index.

Finally, overlap of the bit supports of $n$ and $2n$ is exactly the
condition $n\mathbin{\&}(2n)\ne0$, proving the equivalence with (iii).
\end{proof}

\begin{remark}[The top carry is essential]
For $n=3$, the addition is $11_2+110_2=1001_2$. There are two carries,
not one. The valuation is $1-2+2=1$, in agreement with $a_3=6$.
Omitting the outgoing carry above the leading digit would give the wrong
parity. This boundary condition is explicitly included in every automaton
and program below.
\end{remark}

\section{An independent proof with generating functions}\label{sec:formal}

The ordinary generating function of the sequence can be expressed through
the ternary-tree series. This representation is recorded in A000139
\cite{OEIS139internal}; we derive it to make the second proof independent
of an unproved generating-function identity.

\begin{lemma}\label{lem:ternary}
There is a unique series $T(x)\in1+x\Z[[x]]$ such that
\begin{equation}\label{eq:ternary}
 T(x)=1+xT(x)^3.
\end{equation}
Moreover, with $A(x)=\sum_{n\geq0}a_nx^n$,
\begin{equation}\label{eq:Aternary}
 A(x)=T(x)(3-T(x)).
\end{equation}
In particular, every $a_n$ is an integer.
\end{lemma}
\begin{proof}
Equation \eqref{eq:ternary} determines each nonconstant coefficient from
earlier ones, proving existence, uniqueness, and integrality recursively.
Lagrange inversion, in the special form proved in
\cref{sec:lagrange}, gives, for $k\geq1$,
\[
 [x^n]T(x)^k=\frac{k}{3n+k}\binom{3n+k}{n}.
\]
For $n\geq1$, the coefficients of $T$ and $T^2$ are respectively
\[
 \frac{(3n)!}{n!(2n+1)!},\qquad
 \frac{2(3n+1)!}{n!(2n+2)!}.
\]
Consequently
\[
 [x^n](3T-T^2)
 =\frac{(3n)!}{n!(2n+1)!}
   \left(3-\frac{3n+1}{n+1}\right)
 =\frac{2(3n)!}{(n+1)!(2n+1)!}.
\]
The constant coefficient is $3-1=2$, as required.
\end{proof}

Let
\[
 B(x)=\sum_{\substack{n\geq0\\n\text{ has no binary }11}}x^n
 \quad\text{in }\Z[[x]].
\]
An admissible even integer is $2q$ for an admissible $q$, while an
admissible odd integer is $4q+1$ for an admissible $q$. Hence
\begin{equation}\label{eq:mahler}
 B(x)=B(x^2)+xB(x^4).
\end{equation}
Write bars for reduction of coefficients modulo $2$. Frobenius gives
$\overline B(x^2)=\overline B(x)^2$ and
$\overline B(x^4)=\overline B(x)^4$. Since the constant term of $B$ is
$1$, we may divide the reduced equation by $\overline B$ and obtain
\[
 1=\overline B+x\overline B^3.
\]
Thus $\overline B$ satisfies the same equation and initial condition as
$\overline T$, so $\overline B=\overline T$. Reducing
\eqref{eq:Aternary} gives
\begin{equation}\label{eq:mod2A}
 \overline A=\overline T+\overline T^2
 =\overline B+\overline B^2=x\overline B(x^4).
\end{equation}
The last series has coefficient $1$ exactly at the indices $4q+1$ with
$q$ having no adjacent binary $1$s. These are precisely the odd fibbinary
numbers, proving \cref{thm:parity} a second time.

\begin{corollary}[An algebraic parity series]\label{cor:paritycubic}
The parity series $P(x)=\overline A(x)\in\F_2[[x]]$ satisfies
\begin{equation}\label{eq:paritycubic}
 P(x)=x+x^2P(x)^3.
\end{equation}
It has algebraic degree exactly $3$ over $\F_2(x)$.
\end{corollary}
\begin{proof}
Equation \eqref{eq:mod2A} says $P=x\overline T(x^4)$, which immediately
gives \eqref{eq:paritycubic}. To show that
$x^2Y^3+Y+x$ is irreducible over $\F_2(x)$, it suffices to exclude a
rational root. For a reduced fraction $p/q$ with $p,q\in\F_2[x]$, the
rational-root argument forces $p\mid x$ and $q\mid x^2$. The only
candidates are $1,x,x^{-1},x^{-2}$, none of which is a root.
\end{proof}

\section{Counting odd terms and the limits of a power-law description}\label{sec:oddcount}

\begin{proposition}\label{prop:fibcount}
For every $m\geq0$,
\begin{equation}\label{eq:oddcount}
 \#\{0\leq n<2^m:a_n\text{ is odd}\}=F_m.
\end{equation}
\end{proposition}
\begin{proof}
For $m=0$ the count is $0$, and for $m=1$ it is $1$. When $m\geq2$,
an odd admissible binary word of length $m$, allowing leading zeros,
ends in $01$. Its remaining $m-2$ digits can be any binary word without
$11$. If $b_j$ counts such words of length $j$, splitting according to
the initial block $0$ or $10$ gives $b_j=b_{j-1}+b_{j-2}$, with
$b_0=1$, $b_1=2$. Thus $b_j=F_{j+2}$ and the desired count is $F_m$.
\end{proof}

Put $O(X)=\#\{0\leq n<X:a_n\text{ is odd}\}$ for real $X>0$.
By \eqref{eq:oddcount} and monotonicity between successive powers of $2$,
\begin{equation}\label{eq:oddgrowth}
 O(X)=\Theta(X^\alpha),\qquad \alpha=\log_2\ph<1.
\end{equation}
The odd terms therefore have natural density zero. At dyadic bounds there
is the sharper exact expression
\[
 O(2^m)=\frac{\ph^m-(-\rho)^m}{\sqrt5}.
\]

One should not replace \eqref{eq:oddgrowth} by an unrestricted asymptotic
$O(X)\sim C X^\alpha$ with a single positive constant. All integers in
$[3\cdot2^{m-1},2^{m+1})$ begin with the binary digits $11$, and so none
is an odd index. Consequently
\[
 O(3\cdot2^{m-1})=O(2^{m+1})=F_{m+1}.
\]
The two arguments have a fixed ratio $4/3$, which would contradict the
ratio $(4/3)^\alpha$ predicted by such an asymptotic. The digital
oscillation is a genuine feature, not an error term that disappears at
all scales.

There is also a useful distinction between two kinds of generating
functions. The index series
\[
 H(x)=\sum_{\substack{n\geq0\\a_n\text{ odd}}}x^n\in\Z[[x]]
\]
is not rational over $\Q$. To see this, a rational series with coefficients
in the finite set $\{0,1\}$ satisfies an eventual fixed linear recurrence;
the finite vector of consecutive coefficients has only finitely many
possible values and determines the next vector, so the sequence would
be ultimately periodic. Here there are infinitely many $1$s but also
arbitrarily long zero intervals, which excludes ultimate periodicity.
By contrast, the \emph{binary-length counting series} is rational:
\[
 \sum_{m\geq0}O(2^m)z^m=\frac{z}{1-z-z^2}.
\]
Confusing these two series would conceal an important change of index.

\section{A finite-state model for all valuations}\label{sec:automaton}

\subsection{The three-state carry core}
Read digits of $n$ from least to most significant. A raw carry state
consists of the previous input digit and the incoming carry. There are
four possibilities, but $(1,0)$ and $(0,1)$ have identical transitions
and can be merged. Let
\[
 A=(0,0),\qquad B=\{(1,0),(0,1)\},\qquad C=(1,1).
\]
Weight an outgoing carry by $y$. The transition table is
\begin{center}
\begin{tabular}{c cc c}
\toprule
State & Input $0$ & Input $1$ & Terminal weight\\
\midrule
$A$ & $A$ & $B$ & $1$\\
$B$ & $A$ & $yC$ & $1$\\
$C$ & $yB$ & $yC$ & $y$\\
\bottomrule
\end{tabular}
\end{center}
The notation $yC$ means move to $C$ and multiply the path weight by $y$.
The terminal weight flushes the remaining carry through zero digits.
In particular, a terminal state $C$ creates one more carry and then
moves to $B$, after which the carry vanishes.

The transfer matrix and terminal vector are
\begin{equation}\label{eq:M}
 M(y)=\begin{pmatrix}1&1&0\\1&0&y\\0&y&y\end{pmatrix},
 \qquad f(y)=\begin{pmatrix}1\\1\\y\end{pmatrix}.
\end{equation}
Writing $e_A=(1,0,0)$, the polynomial
\begin{equation}\label{eq:Cm}
 C_m(y)=e_A M(y)^m f(y)
\end{equation}
counts all carry totals $c(n)$ for $0\leq n<2^m$. All length-$m$ words
are read, including those with leading zeros. Such zeros leave the final
weight unchanged.

\subsection{Removing the trailing-run correction without negative weights}
Formula \eqref{eq:carryval} contains $-t(n)$, which could suggest negative
powers of $y$. They are unnecessary. Introduce three prefix states:
$S$ before any digit is read, $U$ after an initial run of exactly one $1$,
and $W$ while an initial run of at least two $1$s is still continuing.
The complete valuation automaton has six states:
\begin{center}
\begin{tabular}{c cc c}
\toprule
State & Input $0$ & Input $1$ & Terminal weight\\
\midrule
$S$ & $yA$ & $U$ & $y$\\
$U$ & $A$ & $W$ & $1$\\
$W$ & $yB$ & $W$ & $y$\\
$A$ & $A$ & $B$ & $1$\\
$B$ & $A$ & $yC$ & $1$\\
$C$ & $yB$ & $yC$ & $y$\\
\bottomrule
\end{tabular}
\end{center}

\begin{proposition}\label{prop:sixstate}
Starting at $S$, the product of the transition weights and terminal weight
for the binary expansion of $n$ is exactly $y^{v(n)}$. This remains true
when arbitrary leading zeros are appended to the input word.
\end{proposition}
\begin{proof}
An even index begins with a $0$. Its valuation is $1+c(n)$, accounted for
by $S\to yA$ and the carry core. An index with $t=1$ begins with the
least-significant-first block $10$, leading from $S$ through $U$ to $A$
with weight $1$; the remaining valuation is precisely the remaining
carry count. The isolated one-digit word has valuation $0$, explaining
the terminal weight at $U$.

For $t\geq2$, the $t$ initial $1$s create $t-1$ carries before the first
zero. In $1-t+c$, their contribution is $1-t+(t-1)=0$, so the loop at $W$
has weight $1$. The terminating zero creates one additional carry and
leaves the raw state $(0,1)$, namely $B$. This gives the transition
$W\to yB$. If the word ends before that zero is read, the mandatory
flush produces exactly this same factor $y$, explaining the terminal
weight at $W$.

The empty word represents $n=0$, so its terminal weight must be $y$.
All other terminal weights are those of the carry core. The rules also
show directly that reading a further zero and then terminating has the
same effect as terminating immediately, in every state.
\end{proof}

This proposition gives a finite combinatorial interpretation with
nonnegative weights for the entire valuation distribution. The extra
prefix states are essential for the factorial denominator; the
three-state core alone counts $c(n)$, not $v(n)$.

\section{The complete dyadic distribution and its rational generating function}\label{sec:dist}

Define
\begin{equation}\label{eq:Vm}
 V_m(y)=\sum_{0\leq n<2^m}y^{v(n)},\qquad
 \mathcal V(z,y)=\sum_{m\geq0}V_m(y)z^m.
\end{equation}
The variable $z$ measures binary length, whereas $y$ measures valuation.

\begin{theorem}[Bivariate distribution]\label{thm:GF}
One has
\begin{equation}\label{eq:mainGF}
 \boxed{\quad
 \mathcal V(z,y)=
 \frac{y+(1-y^2)z-y^3z^2+y(y-1)^2z^3}
 {1-(1+y)z-(y^2-y+1)z^2+y(y+1)z^3}.
 \quad}
\end{equation}
In particular, for $m\geq4$,
\begin{equation}\label{eq:recurrence}
 V_m=(1+y)V_{m-1}+(y^2-y+1)V_{m-2}-y(y+1)V_{m-3},
\end{equation}
with
\begin{equation}\label{eq:seeds}
 V_0=y,\quad V_1=1+y,\quad V_2=1+3y,\quad V_3=2+5y+y^3.
\end{equation}
\end{theorem}
\begin{proof}
Put
\begin{align}
 D(z,y)&=\det(I-zM(y))\notag\\
 &=1-(1+y)z-(y^2-y+1)z^2+y(y+1)z^3.\label{eq:D}
\end{align}
The generating functions for carry-core paths starting at $A$ and $B$
are respectively
\begin{align}
 C(z,y)&=e_A(I-zM)^{-1}f
       =\frac{(1+z)(1-yz)}{D(z,y)},\label{eq:CGF}\\
 R(z,y)&=e_B(I-zM)^{-1}f
       =\frac{1-yz+y^2z(1-z)}{D(z,y)}.\label{eq:RGF}
\end{align}
These expressions can be verified without matrix inversion: if the
three starting-state series are $C,R,Q$, the transition table gives
\[
 C=1+z(C+R),\quad R=1+z(C+yQ),\quad Q=y+zy(R+Q),
\]
and solving these three linear equations gives \eqref{eq:CGF} and
\eqref{eq:RGF}.

Partition all words according to their trailing run. The empty word
contributes $y$. The even words contribute $yzC$. The isolated one-bit
word $1$ contributes $z$, and the words beginning with the
least-significant-first block $10$ contribute $z^2C$. All-ones words of
length at least two contribute $yz^2/(1-z)$. Finally, a run of at least
two $1$s followed by its terminating zero contributes
$yz^3R/(1-z)$. Therefore
\begin{equation}\label{eq:partitionGF}
 \mathcal V=y+yzC+z+z^2C+\frac{yz^2}{1-z}
                         +\frac{yz^3}{1-z}R.
\end{equation}
Substituting \eqref{eq:CGF}--\eqref{eq:RGF} and simplifying yields
\eqref{eq:mainGF}; the apparent factor $1-z$ cancels. Multiplication by
the denominator and extraction of $z^m$ gives \eqref{eq:recurrence} for
$m\geq4$, and direct extraction gives \eqref{eq:seeds}.
\end{proof}

\begin{remark}[The starting index of the recurrence]
The numerator in \eqref{eq:mainGF} has degree three in $z$. Thus the
homogeneous recurrence starts at $m=4$, not at $m=3$. Using only three
initial polynomials and starting too early loses the exceptional prefix
contribution $y(y-1)^2z^3$.
\end{remark}

For illustration, the next polynomials are
\begin{align*}
 V_4&=3+8y+y^2+3y^3+y^4,\\
 V_5&=5+13y+2y^2+7y^3+3y^4+2y^5,\\
 V_6&=8+21y+4y^2+14y^3+8y^4+6y^5+3y^6.
\end{align*}
These show, for example, that valuation $3$ appears before valuation $2$.
The specializations
\begin{equation}\label{eq:specializations}
 \mathcal V(z,0)=\frac{z}{1-z-z^2},\qquad
 \mathcal V(z,1)=\frac1{1-2z}
\end{equation}
recover the Fibonacci odd count and the total count $2^m$.

\section{Every fixed valuation: exact recurrences and asymptotics}\label{sec:fixed}

For a fixed integer $r\geq0$, set
\[
 b_r(m)=[y^r]V_m(y),\qquad H_r(z)=\sum_{m\geq0}b_r(m)z^m.
\]
Write
\begin{equation}\label{eq:d}
 d(z)=1-z-z^2.
\end{equation}
The denominator in \eqref{eq:mainGF} has the particularly useful form
\begin{equation}\label{eq:Dsplit}
 D(z,y)=(1-yz)d(z)-y^2z^2(1-z).
\end{equation}
This explains why the order of the fixed-level recurrence increases
only after every second valuation.

\subsection{A finite closed expression for each level}
Let
\[
 N_0=z,\quad N_1=1+z^3,\quad N_2=-z-2z^3,\quad N_3=z^3-z^2,
\]
so that the numerator of \eqref{eq:mainGF} is $\sum_{j=0}^3N_jy^j$.
For $s\geq0$, define
\begin{equation}\label{eq:J}
 J_s(z)=z^s\sum_{j=0}^{\lfloor s/2\rfloor}
       \binom{s-j}{j}\frac{(1-z)^j}{d(z)^{j+1}},
 \qquad J_s=0\quad(s<0).
\end{equation}
Then
\begin{equation}\label{eq:Hrfinite}
 \boxed{\ H_r(z)=\sum_{j=0}^3N_j(z)J_{r-j}(z).\ }
\end{equation}
Indeed, expand \eqref{eq:Dsplit} as a geometric series:
\[
 \frac1D=\frac1d\sum_{j\geq0}
 \frac{y^{2j}z^{2j}(1-z)^j}{d^j(1-yz)^{j+1}}.
\]
The coefficient of $y^s$ is exactly \eqref{eq:J}, proving
\eqref{eq:Hrfinite}. This is a nonrecursive finite formula for every
fixed-level generating function.

The first four simplify to
\begin{align}
 H_0(z)&=\frac{z}{d},\label{eq:H0}\\
 H_1(z)&=\frac{1+z^2+z^3}{d},\label{eq:H1}\\
 H_2(z)&=\frac{z^4(1-z^2)}{d^2},\label{eq:H2}\\
 H_3(z)&=\frac{z^3(1-z^2)(1+z+z^2)}{d^2}.\label{eq:H3}
\end{align}
For $r\geq4$ they also satisfy the simple recurrence in valuation level
\begin{equation}\label{eq:Hrec}
 H_r=zH_{r-1}+\frac{z^2(1-z)}{d}H_{r-2}.
\end{equation}
All these formulas are identities of formal series; no convergence
assumption is being made.

\begin{theorem}[Minimal eventual recurrences]\label{thm:minrec}
Let $k=\lfloor r/2\rfloor$. In lowest terms over $\Q(z)$, the denominator
of $H_r(z)$ is exactly
\begin{equation}\label{eq:mind}
 (1-z-z^2)^{k+1}.
\end{equation}
Hence $b_r(m)$ has a minimal eventual constant-coefficient recurrence of
order $2k+2$, with characteristic polynomial
\[
 (X^2-X-1)^{k+1}.
\]
\end{theorem}
\begin{proof}
Formula \eqref{eq:Hrfinite} shows that the denominator divides
$d^{k+1}$. To prove that this power cannot be lowered, it is enough to
show a pole of order $k+1$ at the positive root $\rho$ of $d$.
For $r=2k$, only the term $N_0J_{2k}$ contributes to that highest order;
its numerator at $\rho$ is
\[
 \rho^{2k+1}(1-\rho)^k=\rho^{4k+1}>0.
\]
For $r=2k+1$, the two contributing terms are $N_0J_{2k+1}$ and
$N_1J_{2k}$. Their combined numerator at $\rho$ is
\begin{align*}
 &(k+1)\rho^{2k+2}(1-\rho)^k
 +(1+\rho^3)\rho^{2k}(1-\rho)^k\\
 &\hspace{1cm}=\rho^{4k-1}(1+k\rho^3)>0.
\end{align*}
Here $1-\rho=\rho^2$ and $1+\rho=\rho^{-1}$ were used. Since $d$ is
irreducible over $\Q$, this proves \eqref{eq:mind}. A rational generating
function in lowest terms has precisely that denominator as its minimal
eventual recurrence polynomial, with the coefficients reversed.
\end{proof}

\subsection{Explicit leading asymptotics and full expansions}
\begin{theorem}[Fixed-level asymptotics]\label{thm:fixedasymp}
For every fixed $k\geq0$, as $m\to\infty$,
\begin{align}
 b_{2k}(m)&\sim
 \frac{\rho^{3k}}{5^{(k+1)/2}k!}\,m^k\ph^m,\label{eq:evenasymp}\\
 b_{2k+1}(m)&\sim
 \frac{\rho^{3k-2}(1+k\rho^3)}{5^{(k+1)/2}k!}\,m^k\ph^m.
 \label{eq:oddasymp}
\end{align}
More precisely, for each $r$ there are polynomials $P_r,Q_r$ of degree
at most $\lfloor r/2\rfloor$ such that, for all sufficiently large $m$,
\begin{equation}\label{eq:exactBinet}
 b_r(m)=P_r(m)\ph^m+Q_r(m)(-\rho)^m.
\end{equation}
Their coefficients are obtained explicitly by partial fractions of
\eqref{eq:Hrfinite}.
\end{theorem}
\begin{proof}
Factor
\[
 d(z)=(1-z/\rho)(1+\rho z).
\]
At $z=\rho$, the second factor is $1+\rho^2=\sqrt5\,\rho$.
The highest-pole numerators computed in \cref{thm:minrec} therefore give
\begin{align*}
 H_{2k}(z)&\sim
 \frac{\rho^{3k}}{5^{(k+1)/2}}(1-z/\rho)^{-k-1},\\
 H_{2k+1}(z)&\sim
 \frac{\rho^{3k-2}(1+k\rho^3)}{5^{(k+1)/2}}
 (1-z/\rho)^{-k-1}.
\end{align*}
The coefficient of $(1-z/\rho)^{-k-1}$ is
$\binom{m+k}{k}\rho^{-m}\sim m^k\ph^m/k!$, proving the displayed
asymptotics. The only other pole is at $-\ph$. Partial fractions at the
two poles, together with a possible polynomial part of finite support,
give \eqref{eq:exactBinet}. This also provides all lower powers of $m$
in the asymptotic expansion.
\end{proof}

For $k\geq1$, the error after retaining just the leading term is
$O(m^{k-1}\ph^m)$. For $k=0$ there are exact Fibonacci descriptions:
$b_0(m)=F_m$ for all $m$, and $b_1(m)=F_{m+2}$ for $m\geq2$.
The exceptions in the latter statement are $b_1(0)=b_1(1)=1$.

\begin{corollary}[Almost all terms have any prescribed fixed power of $2$]\label{cor:density}
For each fixed $R\geq1$, the set of indices $n$ such that $2^R\mid a_n$
has natural density $1$. More quantitatively, with
$k=\lfloor(R-1)/2\rfloor$,
\[
 \#\{0\leq n<X:2^R\nmid a_n\}
 =\Theta\bigl((\log X)^kX^\alpha\bigr).
\]
\end{corollary}
\begin{proof}
At $X=2^m$, sum the finitely many levels $0\leq r<R$ and use
\cref{thm:fixedasymp}. The largest polynomial degree is $k$, and every
leading constant is positive. Monotonicity between adjacent powers of
$2$ gives the arbitrary-bound estimate. Since $\alpha<1$, its ratio to
$X$ tends to zero.
\end{proof}

These are fixed-$r$ results. They are not uniform estimates when $r$
grows proportionally to $m$. The central regime requires a different
analysis, given in \cref{sec:limit}.

\section{Maximal valuations and the least index of every valuation}\label{sec:extremes}

\begin{theorem}[Sharp dyadic maximum]\label{thm:max}
For every $m\geq3$,
\begin{equation}\label{eq:max}
 \max_{0\leq n<2^m}v(n)=m,
 \qquad \#\{0\leq n<2^m:v(n)=m\}=F_{m-2}.
\end{equation}
Equality $v(n)=m$ occurs exactly for the even indices $n=2q$ such that
$q$ has binary length $m-1$, ends in $11$, and contains no $00$.
The leading digit of $q$ is, as usual, $1$.
\end{theorem}
\begin{proof}
If an integer $q$ has binary length $\ell$, its addition to $2q$ has no
carry in column $0$ and has at most one outgoing carry in each column
$1,\ldots,\ell$. Thus $c(q)\leq\ell$.

For even $n=2q$, shifting does not change the number of carries, so
$v(n)=1+c(q)\leq m$. To have equality, $q$ must have length $\ell=m-1$
and a carry in every column $1,\ldots,\ell$. The first such carry
requires its two lowest digits to be $11$. Once the carry is present,
it survives a column exactly when the two adjacent input digits are not
both $0$. The leading digit $1$ ensures the final outgoing carry.
This proves the stated characterization for even indices.

If $n$ is odd with $t(n)\geq2$, then
$v(n)=1-t(n)+c(n)\leq m-1$. If $t(n)=1$, the two lowest digits are $01$.
No carry can occur in columns $0,1,2$; equivalently, after deleting the
least-significant-first block $10$, the remaining carry computation
starts at $A$. Hence $v(n)\leq m-2$ for $m\geq2$. No odd index is an
extremizer.

It remains to count the described words $q$. Delete their last $1$.
The resulting words have length $m-2$, start and end in $1$, and contain
no $00$. The number of length-$j$ words with these properties is $F_j$:
for $j=1,2$ it is $1$, and splitting after the initial $1$ gives the
Fibonacci recurrence. Therefore there are $F_{m-2}$ extremizers, and
there is at least one for each $m\geq3$.
\end{proof}

For $m=0,1,2$, the maximum is $1$, as the initial polynomials show.
The transition from these exceptional cases to \eqref{eq:max} is why
the least occurrence of valuation $2$ is later than that of valuation
$3$.

\begin{theorem}[Least positive indices]\label{thm:least}
Let $\ell_r$ be the least positive integer $n$ such that $v(n)=r$.
Then
\begin{equation}\label{eq:least}
 \ell_0=1,\qquad\ell_1=2,\qquad\ell_2=13,
\end{equation}
and for all $r\geq3$,
\begin{equation}\label{eq:leastgeneral}
 \boxed{\ \ell_r=\frac{2^{r+1}+6+4(-1)^r}{3}.\ }
\end{equation}
\end{theorem}
\begin{proof}
The valuations at $n=1,\ldots,13$ are
\[
 0,1,1,1,0,3,1,1,0,1,3,3,2,
\]
which proves \eqref{eq:least}. For $r\geq3$, \cref{thm:max} shows that
no index of fewer than $r$ binary digits can have valuation $r$.
Among indices of length $r$, the extremal-word characterization applies.
The smallest eligible binary word is obtained lexicographically by
choosing $0$ whenever it does not create $00$ or violate the final $11$
condition on $q$. Consequently the binary word of $\ell_r$ is
\[
 \begin{cases}
 (10)^{(r-3)/2}110,&r\text{ odd},\\
 (10)^{(r-4)/2}1110,&r\text{ even}.
 \end{cases}
\]
Summing the corresponding geometric progressions gives
$(2^{r+1}+2)/3$ in the odd case and $(2^{r+1}+10)/3$ in the even case,
which is \eqref{eq:leastgeneral}.
\end{proof}

The beginning of this least-index sequence is
\[
 1,2,13,6,14,22,46,86,174,342,686,1366,\ldots.
\]
Its tail is a linear combination of $2^r$, $1$, and $(-1)^r$.
Accordingly, for $r\geq6$,
\[
 \ell_r=2\ell_{r-1}+\ell_{r-2}-2\ell_{r-3}.
\]
Here ``least'' refers to a positive index. Allowing the index $0$ would
replace $\ell_1=2$ by $0$ and leave the other values unchanged.

\section{Exact moments and Gaussian limit laws}\label{sec:limit}

Let $N_m$ be uniformly distributed on $\{0,1,\ldots,2^m-1\}$ and put
$X_m=v(N_m)$. Its probability generating function is
\begin{equation}\label{eq:pgf}
 \E[y^{X_m}]=2^{-m}V_m(y).
\end{equation}
This is a random \emph{index} model. It is not a model that first chooses
a permutation uniformly among all objects of several different sizes.

\subsection{Exact mean and variance}
\begin{theorem}\label{thm:moments}
For all $m\geq0$,
\begin{equation}\label{eq:mean}
 \E X_m=\frac m2-\frac23+\frac{9+(-1)^m}{6\,2^m}.
\end{equation}
For $m\geq1$,
\begin{equation}\label{eq:variance}
 \Var X_m=\frac{11m}{12}-\frac{23}{9}
 +\frac{27-(-1)^m}{6\,2^m}
 -\frac{(9+(-1)^m)^2}{36\,4^m}.
\end{equation}
For $m=0$, $X_0=1$ deterministically and its variance is $0$.
\end{theorem}
\begin{proof}
Differentiate \eqref{eq:mainGF} with respect to $y$ and set $y=1$.
The resulting first-moment series has the partial-fraction form
\begin{equation}\label{eq:firstpartial}
 \sum_{m\geq0}V_m'(1)z^m
 =-\frac{7}{6(1-2z)}+\frac1{2(1-2z)^2}
   +\frac1{6(1+z)}+\frac3{2(1-z)}.
\end{equation}
Thus
\[
 V_m'(1)=2^m\left(\frac m2-\frac23\right)
              +\frac32+\frac{(-1)^m}{6},
\]
which proves \eqref{eq:mean}.

Differentiating twice gives the factorial second moment. The full
partial-fraction identity is recorded in \cref{sec:partials}; its
coefficient is
\begin{align}\label{eq:factorialmoment}
 V_m''(1)={}&2^m\left(\frac{m^2}4-\frac m4-\frac{13}9\right)
 +(-1)^m\left(\frac m6-\frac59\right)\notag\\
 &+\frac{3m}{2}+1+\ind_{m=0}.
\end{align}
Since $\E X_m^2=2^{-m}(V_m''(1)+V_m'(1))$, subtracting the square of
\eqref{eq:mean} and simplifying gives \eqref{eq:variance} for $m\geq1$.
The indicator in \eqref{eq:factorialmoment} accounts for the exceptional
constant coefficient when $m=0$.
\end{proof}

In particular,
\begin{equation}\label{eq:momentasymp}
 \E X_m=\frac m2-\frac23+O(2^{-m}),\qquad
 \Var X_m=\frac{11m}{12}-\frac{23}{9}+O(2^{-m}).
\end{equation}
The variance is substantially larger than $m/4$, the variance of a sum
of $m$ independent fair binary digits. Carry propagation creates
correlations; merely counting independent local patterns would miss
this effect.

\subsection{The dominant eigenvalue and its symmetry}
The characteristic polynomial of \eqref{eq:M} is
\begin{equation}\label{eq:char}
 \chi(\lambda,y)=\lambda^3-(1+y)\lambda^2
 -(y^2-y+1)\lambda+y(y+1).
\end{equation}
At $y=1$ it factors as $(\lambda-2)(\lambda-1)(\lambda+1)$.
Let $\lambda(y)$ be the analytic root with $\lambda(1)=2$.
For real $y>0$ near $1$, it is the Perron eigenvalue of the nonnegative,
irreducible, aperiodic matrix $M(y)$.

The simple pole corresponding to this root in \eqref{eq:mainGF} gives
analytic functions $h(y)$ and an $\eta<2$, on a sufficiently small complex
neighborhood of $1$, such that
\begin{equation}\label{eq:quasipower}
 V_m(y)=h(y)\lambda(y)^m+O(\eta^m),\qquad h(1)=1.
\end{equation}
The error is uniform in a smaller closed neighborhood. This follows
either from the spectral projections of $M(y)$ or directly by partial
fractions in $z$, since the three roots remain separated there. The
polynomial part in $z$ affects only finitely many initial $m$.

Let $J$ be the permutation matrix swapping states $A$ and $C$. A direct
calculation gives
\[
 JM(y)J=yM(1/y).
\]
Perron eigenvalues for positive real $y$, followed by analytic
continuation near $1$, therefore satisfy
\begin{equation}\label{eq:symmetry}
 \lambda(y)=y\lambda(1/y).
\end{equation}
Put
\[
 K(s)=\log\frac{\lambda(e^s)}2.
\]
Equation \eqref{eq:symmetry} says that $K(s)-s/2$ is even. In particular,
all its odd Taylor coefficients beyond the first vanish.
Implicit differentiation of \eqref{eq:char} gives
\[
 \lambda'(1)=1,\qquad\lambda''(1)=\frac43,
\]
and hence
\begin{equation}\label{eq:Ktaylor}
 K'(0)=\frac12,\qquad K''(0)=\frac{11}{12}.
\end{equation}
A further algebraic expansion, reproducible by the symbolic verification
script, gives
\begin{equation}\label{eq:Kexpansion}
 K(s)=\frac s2+\frac{11s^2}{24}-\frac{569s^4}{1728}+O(s^6).
\end{equation}
For example, one can first solve \eqref{eq:char} in powers of $u=y-1$:
\[
 \lambda(1+u)=2+u+\frac23u^2-\frac13u^3-\frac7{27}u^4+O(u^5),
\]
and then substitute $u=e^s-1$. Symmetry supplies the vanishing fifth
coefficient. The usual analytic generating-function framework for such
limit laws is described in \cite{FS}; the needed argument is included
next rather than invoked as a black box.

\subsection{Central and local limit theorems}
\begin{theorem}\label{thm:CLT}
As $m\to\infty$,
\begin{equation}\label{eq:CLT}
 \frac{X_m-m/2}{\sqrt{11m/12}}
 \ \xrightarrow{\ d\ }\ \mathcal N(0,1).
\end{equation}
The following stronger local statement holds uniformly for integers $r$:
\begin{equation}\label{eq:LCLT}
 \sup_{r\in\Z}\left|
 \sqrt{\frac{11m}{12}}\,\Prob(X_m=r)
 -\frac1{\sqrt{2\pi}}
  \exp\!\left(-\frac{(r-m/2)^2}{2(11m/12)}\right)
 \right|\longrightarrow0.
\end{equation}
One may replace $m/2$ by the exact mean in either statement.
\end{theorem}
\begin{proof}
Write $\sigma^2=11/12$ and $\mu=1/2$. For fixed real $t$, substitute
$y=\exp(it/(\sigma\sqrt m))$ into \eqref{eq:quasipower}. After multiplying
by $\exp(-it\mu\sqrt m/\sigma)$, \eqref{eq:Ktaylor} yields the limit
$\exp(-t^2/2)$ for the characteristic function of the normalized
variable. The prefactor $h(y)$ tends to $1$, and the error is
exponentially small. The continuity theorem for characteristic functions
proves \eqref{eq:CLT}.

We verify the extra spectral condition for a local theorem. For $|u|=1$,
entrywise absolute values give $|M(u)|=M(1)$, so every eigenvalue has
modulus at most $2$. Suppose $M(u)w=\gamma w$ with $|\gamma|=2$.
Taking absolute values gives $2|w|\leq M(1)|w|$. Multiplication by the
positive left Perron vector forces equality in every coordinate. By
irreducibility, $|w|$ is a positive Perron vector, and equality also
holds in every triangle inequality used above. The self-loop at $A$,
whose weight is $1$, then forces $\gamma/2=1$. The self-loop at $C$,
whose weight is $u$, forces $u=1$. Therefore
\begin{equation}\label{eq:strictgap}
 \operatorname{spr}(M(e^{it}))<2
 \quad\text{for }0<|t|\leq\pi.
\end{equation}
On a closed arc bounded away from $0$, compactness makes this a uniform
strict gap. The recurrence \eqref{eq:recurrence}, whose characteristic
roots are those of $M$, consequently bounds $V_m(e^{it})/2^m$
uniformly by $C\theta^m$ there, for some $\theta<1$. A resolvent contour
with radius strictly between the maximum root modulus and $2$ gives
this bound even at a repeated subdominant eigenvalue.

Near $t=0$, \eqref{eq:Ktaylor} and analyticity give constants $c,\delta>0$
such that
\[
 \left|\frac{\lambda(e^{it})}{2}\right|^m\leq e^{-cm t^2}
 \qquad(|t|\leq\delta).
\]
Fourier inversion gives
\[
 \Prob(X_m=r)=\frac1{2\pi}\int_{-\pi}^{\pi}
                 2^{-m}V_m(e^{it})e^{-irt}\dd t.
\]
Set $s=\sigma\sqrt m\,t$ and remove the phase corresponding to $\mu m$.
On each bounded $s$-interval, \eqref{eq:quasipower} gives convergence of
the integrand to $e^{-s^2/2}$. The preceding Gaussian bound dominates the
part with $|t|\leq\delta$, and \eqref{eq:strictgap} makes the remaining
part negligible even after multiplication by $\sqrt m$. Thus the
rescaled characteristic functions converge in $L^1$ to the Gaussian
characteristic function, with the finite integration interval extended
by zero. Fourier inversion and the bound of the difference of transforms
by the $L^1$ difference prove \eqref{eq:LCLT}, uniformly in $r$.

Finally, \eqref{eq:mean} differs from $\mu m$ by a bounded amount. Dividing
that difference by $\sigma\sqrt m$ gives zero in the limit, and the
Gaussian density is uniformly continuous. This justifies either choice
of center.
\end{proof}

The local theorem describes the central distribution on the scale
$\sqrt m$. It does not provide relative errors in the extreme tails.
In particular, the rare fixed-valuation regime is governed by
\cref{thm:fixedasymp}, not by substituting a fixed $r$ into a Gaussian
approximation and expecting relative accuracy.

\section{Algorithms, independent checks, and reproducibility}\label{sec:algorithms}

\subsection{Evaluating a single index}
The exact valuation requires no construction of $a_n$:
\begin{lstlisting}[language=Python]
def valuation(n: int) -> int:
    if not isinstance(n, int) or isinstance(n, bool) or n < 0:
        raise ValueError("n must be a nonnegative integer")
    return n.bit_count() + (n + 1).bit_count() - (3*n).bit_count()
\end{lstlisting}
The parity criterion is the still simpler predicate
\begin{lstlisting}[language=Python]
(n & 1) != 0 and (n & (n << 1)) == 0
\end{lstlisting}
These operate on the digits of $n$, whereas $a_n$ itself has a number of
digits proportional to $n$. The six-state automaton supplies an
alternative digit-by-digit implementation, independent of the
population-count expression.

\subsection{The complete distribution below an arbitrary bound}
The recurrence \eqref{eq:recurrence} computes all $V_0,\ldots,V_m$ with
$O(m^2)$ integer coefficient updates. An arbitrary interval $[0,N)$ can
also be handled without enumeration. Process candidate digits and the
digits of $N$ from least to most significant, maintaining the six-state
automaton and a comparison state in $\{-1,0,1\}$.

After the lowest $j$ digits have been read, the comparison state is the
comparison of those two low-digit integers. At the next position, an
equal pair of bits leaves the comparison unchanged; an unequal pair
\emph{overrides} it, because that new bit is more significant than every
previously read bit. After all positions, retain only comparison $-1$.
This gives an 18-state digit dynamic program. Attach a polynomial in
$y$ to each state and apply the weights from \cref{sec:automaton}.
The terminal weights complete the valuation.

For $L=\bits(N)$ this uses $O(L^2)$ integer coefficient updates and stores
$O(L)$ integer coefficients across a constant number of states. These
are arithmetic-operation counts, not unit-cost bit-complexity claims:
the coefficients themselves may have $O(L)$ bits. The implementation
checks that the final sum of histogram entries equals $N$ exactly.
The case $N=0$ is handled separately by the zero histogram.

\subsection{What was actually executed}
The archive contains a standard-library implementation
\texttt{code/carries.py}, a standard-library verification driver
\texttt{code/verify.py}, and an optional symbolic verifier using SymPy.
The full exact test run recorded the following checks:
\begin{center}
\begin{tabular}{p{0.57\linewidth}p{0.30\linewidth}}
\toprule
Check & Executed range\\
\midrule
Direct factorial quotient against digit valuation, carry count,
automaton valuation, and parity & $0\leq n\leq2000$\\
Digit valuation and parity; histograms against the recurrence
& $0\leq n<2^{20}$\\
Independent column-carry count and six-state automaton against digit sums
& $0\leq n<2^{18}$\\
Every arbitrary-bound digit-DP histogram against enumeration
& $0\leq N\leq2048$\\
Digit DP against dyadic recurrence; exact moments, parity counts,
and extremal coefficients & $0\leq m\leq256$\\
Least positive index, by exhaustive search & $0\leq r\leq18$\\
Formula witnesses of valuation $r$, without exhaustive minimality testing
& $3\leq r\leq1000$\\
Non-dyadic large-bound digit DP and exact total & $N=10^{100}$\\
\bottomrule
\end{tabular}
\end{center}
All these checks passed. The executed environment used Python 3.13.5;
the full standard-library verification run took about 3.3 seconds in
this session. Timing is environment-specific and is not a performance
guarantee. The exact report, including the measured duration, is in
\texttt{data/verification.json}.

These comparisons deliberately distinguish levels of independence.
The million-index test does not compute a million factorial quotients;
it compares the proved digit formula with the binary predicate and the
independently implemented distribution recurrence. Direct quotient
checks stop at $2000$. Likewise, evaluating the formula for $\ell_{1000}$
checks that its valuation is $1000$, but does not exhaust all smaller
integers. Minimality at every level rests on \cref{thm:least}.

The symbolic verifier checks the rational identities, fixed-level
formulas, moment decompositions, and the eigenvalue expansion. It is a
useful algebraic cross-check, not a machine-checked proof of the
Perron--Frobenius or Fourier arguments. No Lean formalization or external
peer review is claimed.

\Needspace{15\baselineskip}
\subsection{Reproduction commands}
From the extracted archive directory, run
\begin{lstlisting}
python code/verify.py
python code/carries.py --n 13
python code/carries.py --dyadic 100
python code/carries.py --bound 1000000000000000000000000
python code/carries.py --least 1000
python code/verify_symbolic.py
pdflatex -interaction=nonstopmode -halt-on-error research.tex
pdflatex -interaction=nonstopmode -halt-on-error research.tex
\end{lstlisting}
The symbolic script is optional and requires SymPy. The main program and
exact verification driver require no third-party packages and make no
network calls. Generated CSV files contain the first $201$ sequence
terms, every dyadic histogram through $m=256$, and least indices through
valuation $64$. The large-bound JSON contains the complete valuation
histogram for $0\leq n<10^{100}$.

\section{Conclusion, provenance, and remaining directions}\label{sec:conclusion}

The selected parity assertion is established by \cref{thm:parity}, and
\cref{sec:formal} gives an independent proof in formal power series.
Its mathematical explanation is especially short: factorial valuations
are digit sums, the sum $n+2n$ is controlled by adjacent digits, and the
factor $n+1$ forces careful treatment of the trailing run of $1$s.

The same analysis gives substantially more than parity. The full dyadic
valuation distribution has the rational generating function
\eqref{eq:mainGF}. The cubic denominator governs a recurrence in binary
length, fixed valuation levels have exact denominators
$d^{\lfloor r/2\rfloor+1}$, and the dominant poles give their asymptotic
expansions. A different dominant-eigenvalue calculation gives the exact
moment constants and the Gaussian central regime. Extremal words give
both the Fibonacci number of maximizers and the explicit least-index
formula. The algorithms directly implement these finite-state proofs.

The status statement should remain precise. The OEIS page provides a
recent, explicitly conjectural target; the article proves that target.
Targeted searches for A000139 parity and valuation, fibbinary indices,
and parity of two-stack-sortable permutations did not locate a prior
resolution in the material retrieved. This is not an exhaustive
bibliographic certification. The classical enumeration formula, the
ternary-tree generating function, and the fibbinary terminology are not
claimed as new. Nor is the general method of using finite matrices to
study digital statistics new. The specific identities proved here are
presented with explicit derivations so that their correctness can be
assessed independently of any priority claim.

Several natural extensions remain outside the scope of the proofs.
One is a parity-reversing involution directly on two-stack-sortable
permutations whose fixed objects explain the fibbinary indices. Another
is the joint distribution of valuations at different primes, where the
binary shift structure no longer supplies this particular three-state
core. A third is a uniform saddle-point analysis when $r/m$ stays away
from both $0$ and the central value $1/2$, connecting the fixed-level
and central regimes through explicit large-deviation asymptotics.
Finally, the full normalized counting oscillation for non-dyadic bounds
can be studied through the digit dynamic program rather than reduced
to upper and lower bounds. No solution of these additional questions is
asserted here.

\appendix
\section{The special Lagrange inversion identity used above}\label{sec:lagrange}

We give the formal calculation to keep \cref{lem:ternary} self-contained.
Let $U=x\Phi(U)$ with $\Phi(0)=1$. For $n\geq1$ and a formal power series
$G$, the coefficient formula is
\begin{equation}\label{eq:lagrange}
 [x^n]G(U(x))=\frac1n[u^{n-1}]G'(u)\Phi(u)^n.
\end{equation}
Indeed, using formal residues and substituting $x=u/\Phi(u)$ gives
\begin{align*}
 [x^n]G(U(x))
 &=\operatorname{Res}_{u=0}
 G(u)\left(\frac{\Phi(u)}u\right)^{n+1}
 \frac{\Phi(u)-u\Phi'(u)}{\Phi(u)^2}\dd u\\
 &=\operatorname{Res}_{u=0}
 G(u)\left(\frac{\Phi(u)^n}{u^{n+1}}
 -\frac{\Phi(u)^{n-1}\Phi'(u)}{u^n}\right)\dd u.
\end{align*}
The expression in parentheses is
$-\frac1n\frac{\dd}{\dd u}(\Phi(u)^n/u^n)$. Formal integration by parts,
whose validity is the fact that a formal derivative has zero residue,
proves \eqref{eq:lagrange}. The residue substitution is valid because
$u/\Phi(u)=u+O(u^2)$ has a compositional inverse; checking monomials gives
the usual change-of-variable rule.

For $T=1+U$ and $\Phi(u)=(1+u)^3$, take $G(u)=(1+u)^k$. Then
\[
 [x^n]T^k=\frac{k}{n}\binom{3n+k-1}{n-1}
 =\frac{k}{3n+k}\binom{3n+k}{n}.
\]
For $n=0$ both sides are $1$, so the same expression also holds there.

\section{Second-moment algebra and a compact symbolic certificate}\label{sec:partials}

The second derivative of \eqref{eq:mainGF} at $y=1$ is
\begin{align}\label{eq:secondpartial}
 \sum_{m\geq0}V_m''(1)z^m
 ={}&1-\frac{17}{18(1-2z)}-\frac1{(1-2z)^2}
       +\frac1{2(1-2z)^3}\notag\\
 &-\frac{13}{18(1+z)}+\frac1{6(1+z)^2}
       -\frac1{2(1-z)}+\frac3{2(1-z)^2}.
\end{align}
Extracting coefficients using
\[
 [z^m](1-az)^{-j}=\binom{m+j-1}{j-1}a^m
\]
gives \eqref{eq:factorialmoment}. This also makes the isolated constant
$1$ on the right, and hence the $m=0$ correction, explicit.

A compact exact certificate for the principal generating-function
identity is the polynomial equality
\[
 D(z,y)\!\left(y+yzC+z+z^2C+\frac{yz^2}{1-z}
 +\frac{yz^3R}{1-z}\right)
 =y+(1-y^2)z-y^3z^2+y(y-1)^2z^3,
\]
with $C,R$ as in \eqref{eq:CGF}--\eqref{eq:RGF}. Multiplication by $1-z$
turns it into a literal identity in $\Z[z,y]$. The included symbolic
script checks this identity after exact cancellation, along with the
fixed-level formulas and moment identities. Thus the displayed rational
function is not inferred by fitting a finite initial segment.

\section{A small exact table}\label{sec:smalltable}

The following table records $b_r(m)$ for the first few binary lengths.
Blank entries are zero. Each row sums to $2^m$; its first entry is $F_m$.
\begin{center}
\setlength{\tabcolsep}{7pt}
\begin{tabular}{c rrrrrrrrr}
\toprule
$m\backslash r$&0&1&2&3&4&5&6&7&8\\
\midrule
0&0&1&&&&&&&\\
1&1&1&&&&&&&\\
2&1&3&&&&&&&\\
3&2&5&0&1&&&&&\\
4&3&8&1&3&1&&&&\\
5&5&13&2&7&3&2&&&\\
6&8&21&4&14&8&6&3&&\\
7&13&34&8&27&16&16&9&5&\\
8&21&55&15&51&32&34&25&15&8\\
\bottomrule
\end{tabular}
\end{center}

The first valuations at individual indices are
\[
\begin{array}{c|rrrrrrrrrrrrrrrr}
 n&0&1&2&3&4&5&6&7&8&9&10&11&12&13&14&15\\\hline
 v(n)&1&0&1&1&1&0&3&1&1&0&1&3&3&2&4&1
\end{array}
\]
These tables are for orientation, not substitutes for the all-index
arguments. The CSV artifacts provide larger exact tables.

\begin{thebibliography}{9}
\bibitem{OEIS139}
The OEIS Foundation Inc.,
\emph{A000139: $a(n)=2(3n)!/((2n+1)!(n+1)!)$}.
Includes Peter Bala's parity conjecture dated July 24, 2025.
Accessed September 19, 2026.
\url{https://oeis.org/A000139}.

\bibitem{OEIS139internal}
The OEIS Foundation Inc., \emph{A000139, internal-format entry}.
Inspected revision 297, May 12, 2026; accessed September 19, 2026.
\url{https://oeis.org/A000139/internal}.

\bibitem{OEIS22341}
The OEIS Foundation Inc., \emph{A022341: the odd fibbinary numbers}.
Accessed September 19, 2026.
\url{https://oeis.org/A022341}.

\bibitem{LSW}
L. Lindroos, A. Sills, and H. Wang,
\emph{Odd fibbinary numbers and the golden ratio},
The Fibonacci Quarterly \textbf{52} (2014), no.~1, 61--65.
\url{https://www.fq.math.ca/Papers1/52-1/LindroosSillsWang.pdf}.

\bibitem{Fang}
W. Fang, \emph{Fighting fish and two-stack sortable permutations},
S\'eminaire Lotharingien de Combinatoire \textbf{80B} (2018),
Article 7, 12 pp.
\url{https://www.mat.univie.ac.at/~slc/wpapers/FPSAC2018/07-Fang.pdf}.
Full-paper preprint: \url{https://arxiv.org/abs/1711.05713}.

\bibitem{FS}
P. Flajolet and R. Sedgewick, \emph{Analytic Combinatorics},
Cambridge University Press, 2009. See the chapters on rational
functions and limit laws. Authors' book page:
\url{https://algo.inria.fr/flajolet/Publications/books.html}.
\end{thebibliography}
\end{document}
